\documentclass[11pt,a4paper]{article}\usepackage[margin=25mm]{geometry}\usepackage{fancyhdr}\usepackage{setspace}\usepackage{titlesec}\usepackage{xcolor}
\usepackage{amsmath,amssymb,graphicx,booktabs,url,array}
\usepackage[T1]{fontenc}
\usepackage{microtype}
\usepackage[hidelinks]{hyperref}
\setlength{\emergencystretch}{1em}
\renewcommand{\arraystretch}{1.15}
\definecolor{reportnavy}{HTML}{123747}
\hypersetup{pdftitle={Molecular Toxicity Prediction: Project Report}}
\pagestyle{fancy}
\fancyhf{}
\fancyhead[L]{\small\textsf{MOLECULAR TOXICITY PREDICTION}}
\fancyhead[R]{\small\textsf{CSAI PROJECT REPORT}}
\fancyfoot[L]{\small Machine Learning and Deep Learning}
\fancyfoot[R]{\small Page \thepage}
\setlength{\headheight}{15pt}
\renewcommand{\headrulewidth}{0.4pt}
\renewcommand{\footrulewidth}{0.4pt}
\titleformat{\section}{\Large\bfseries\color{reportnavy}}{\thesection}{0.8em}{}
\titleformat{\subsection}{\large\bfseries\color{reportnavy}}{\thesubsection}{0.8em}{}
\setstretch{1.12}
\setlength{\parindent}{0pt}
\setlength{\parskip}{6pt}
\begin{document}
\begin{titlepage}
\centering
\vspace*{25mm}
{\large\textsf{CSAI STUDENT PROJECT}\par}
\vspace{18mm}
{\Huge\bfseries\color{reportnavy} Molecular Toxicity\\[6pt] Prediction\par}
\vspace{12mm}
{\Large Machine Learning, Multitask Neural Networks\\[5pt] and Hybrid Random Forests\par}
\vspace{15mm}
{\large Formal Project Report\par}
{\large Phase 1 and Phase 2\par}
\vfill
\begin{tabular}{ll}
\toprule
Team member & Name and roll number\\\midrule
1 & To be completed\\
2 & To be completed\\
3 & To be completed\\
4 & To be completed\\
5 & To be completed\\\bottomrule
\end{tabular}
\vspace{15mm}

{\large November 2026\par}
\end{titlepage}
\pagenumbering{roman}
\tableofcontents
\newpage
\pagenumbering{arabic}
\begin{abstract}
We compare classical machine learning, a multitask neural network, and a hybrid neural-feature Random Forest for 12 Tox21 assay endpoints. The pipeline uses 7,823 valid molecules and a frozen scaffold partition of 6,259 training, 782 validation and 782 test molecules. Inputs combine 2,048 Morgan fingerprint bits with five standardized descriptors. Unmeasured assay labels are excluded from optimization and evaluation. The historical standalone Random Forest reports macro ROC-AUC 0.7541, compared with 0.7179 in the original neural evaluation and 0.6753 for the hybrid. A separate three-seed multitask experiment reports 0.7472 with standard deviation 0.0065. A separately verified inference bundle achieves ROC-AUC 0.7498 (RF), 0.7318 (DL) and 0.7182 (hybrid). Both evaluations favor the standalone forest; hybrid representation learning does not yield a measured improvement. A Streamlit demonstration applies the saved models to new SMILES inputs and exports endpoint-level scores.
\end{abstract}
\noindent\textbf{Keywords:} Tox21; multitask learning; molecular fingerprints; scaffold split; hybrid model.\medskip

\section{Problem and Scope}
The objective is to estimate binary biological activity across seven nuclear receptor and five stress response assays from molecular structure. The 12 outputs are not mutually exclusive. Assay activity is narrower than human toxicity: an inactive prediction does not establish chemical safety. Phase 1 implemented Logistic Regression, Decision Tree and Random Forest baselines. Phase 2 adds shared neural representations, a hybrid classifier, controlled ablations and saved-model inference.

\subsection{Research Questions and Contributions}
The study asks three questions: (1) does a shared neural representation improve prediction across sparsely labelled assays; (2) does a forest benefit from learned molecular embeddings; and (3) can the evaluated models be reproduced at inference time? The contributions are a common data contract, a masked multitask implementation, controlled representation comparisons, and an inference audit that exposed an artifact inconsistency. The hybrid is evaluated as an empirical hypothesis rather than assumed to outperform its components.

\subsection{Project Objectives and Scope}
The complete workflow comprises molecular data validation, exploratory analysis, feature engineering, classical baselines, deep learning, hybrid experiments, technical validation, and a unified inference dashboard. The scope is prediction of assay activity from two-dimensional structure. Clinical toxicity, dose response and regulatory safety assessment are outside the evaluated task.

\section{Related Work}
Multitask networks have been studied for drug discovery, providing a rationale for sharing representations across assays~\cite{ramsundar}. MoleculeNet provides a molecular machine-learning benchmark context and motivates careful split and metric selection~\cite{moleculenet}. Morgan fingerprints encode local molecular neighborhoods~\cite{morgan}, while Random Forests aggregate randomized decision trees~\cite{breiman}. Our approach uses these established techniques on the team's frozen dataset, rather than claiming a new neural architecture. The detailed literature comparison is committed under Member 1's literature directory.


\subsection{How the Literature Informs the Design}
The fingerprint approach provides a fixed-dimensional representation that can be reused by both classical and neural methods. This makes comparisons easier to interpret because a change in predictive model need not also change the input representation. Random Forest provides a strong nonlinear reference with limited preprocessing demands. Multitask learning introduces a different inductive bias: assays share internal representations while retaining separate outputs. Sharing may help when useful structural patterns recur across endpoints, but it can also create negative transfer when endpoints require conflicting features.

The hybrid experiment tests a specific extension of these ideas. A supervised neural encoder learns a compact representation from all assays, while downstream forests make endpoint-specific decisions. This rationale does not guarantee an improvement: a representation optimized for neural output heads may discard detail that a forest would otherwise use. The experimental design therefore retains the original-feature RF and standalone neural model as essential controls.


\section{Dataset and Leakage Controls}
Eight unparseable structures are excluded from the original 7,831 records, leaving 7,823 valid compounds. Approximately 17.06\% of assay entries are unmeasured. The authoritative Phase 2 split is Member 1's frozen split B, matching \texttt{labels\_and\_splits.parquet}. It contains 6,259/782/782 training/validation/test compounds. The older Phase 1 split manifest is a different partition and must not be mixed with these artifacts.

Molecules are aligned explicitly by identifier before consuming features. Identical split membership is retained across model comparisons. Descriptor scaling uses training statistics only. Hyperparameter and checkpoint selection use validation data; the test set is reserved for evaluation. Inference does not read labels or fit preprocessing.

\begin{table}[t]
\caption{Dataset partitions and their permitted uses}
\label{tab:data}\centering
\begin{tabular}{lrl}\toprule
Partition & Molecules & Purpose\\\midrule
Training & 6,259 & Fit models and scaler\\
Validation & 782 & Select settings and checkpoints\\
Test & 782 & Report final performance\\\midrule
Total & 7,823 & 12 assay outputs per molecule\\\bottomrule
\end{tabular}
\end{table}


\subsection{Label Structure and Assay Coverage}
Each molecule has up to 12 labels, with 1 indicating active, 0 indicating inactive, and a missing entry indicating that the corresponding assay outcome is unavailable. There are $7,823\times12=93,876$ possible molecule--assay entries. The data audit records 16,012 missing entries, leaving 77,864 observed labels. Consequently, a compound can contribute to one endpoint's training loss while being excluded from another endpoint's loss in the same minibatch.

\begin{table}[htbp]
\centering\caption{Assay groups represented in the prediction task}
\begin{tabular}{p{0.22\textwidth}p{0.67\textwidth}}\toprule
Group & Endpoints\\\midrule
Nuclear receptors & NR-AR, NR-AR-LBD, NR-AhR, NR-Aromatase, NR-ER, NR-ER-LBD, NR-PPAR-gamma\\
Stress response & SR-ARE, SR-ATAD5, SR-HSE, SR-MMP, SR-p53\\\bottomrule
\end{tabular}
\end{table}

Missingness and imbalance are different problems. Masking addresses unavailable outcomes; positive-class weighting addresses unequal active and inactive counts among observed outcomes. Neither procedure creates new experimental evidence. In particular, imputation of missing outcomes as inactive would increase the apparent negative population and change the supervised task.

\subsection{Scaffold Policy and Distribution Shift}
The frozen partition groups molecules using their structural scaffolds to reduce overlap of molecular cores between training and evaluation. The audit records zero scaffold and molecule-identifier overlap across the three partitions. It also records an important boundary: all 1,775 acyclic molecules occur in training, and none occur in validation or test. The reported evaluation therefore concerns the held-out scaffold population and does not directly estimate performance on unseen acyclic molecules.

A random split can place structurally similar compounds in both training and test data. This can make prediction easier than generalization to new scaffold families. A scaffold split is a useful stress test, but it is not a complete simulation of prospective deployment: assay prevalence, chemical coverage and collection practices can also change. Comparisons between split policies should therefore report the population and prevalence of each evaluation set.

\subsection{Integrity Checks Before Model Fitting}
The pipeline checks that molecule identifiers align across metadata, fingerprints and descriptors. Counts must satisfy observed labels equals active plus inactive labels, and total entries equals observed plus missing entries. The committed audit passes these checks across all endpoints and partitions. All evaluation endpoints contain both classes; the smallest positive counts are 24 in validation and 21 in test. These counts permit ROC-AUC computation, while also indicating why rare-endpoint estimates can be variable.

The scaler is fitted on 6,259 training molecules. Feature matrices and split metadata are saved as reusable artifacts so downstream models share the same data contract. Any future regeneration must preserve molecule ordering or explicitly realign by identifier before fitting models.


\section{Representations and Models}
\subsection{Molecular Features}
Each molecule is represented by a radius-2, 2,048-bit Morgan fingerprint concatenated with molecular weight, topological polar surface area, MolLogP, hydrogen-bond donor count and hydrogen-bond acceptor count. Continuous descriptors are standardized using committed training means and standard deviations. Their positional order is preserved during inference.


\subsection{Feature Construction and Interpretation}
Morgan fingerprints summarize local atom environments using iterative neighborhood expansion. Radius 2 captures neighborhoods extending up to two bonds from an atom; hashing maps the resulting identifiers into 2,048 binary positions. A bit represents the presence of at least one encoded environment rather than a unique, directly interpretable chemical group. Hash collisions and the fixed radius limit the information available to downstream models.

The five continuous descriptors complement the local fingerprint with broad molecular properties. Molecular weight summarizes size, TPSA summarizes polar surface contribution, MolLogP estimates lipophilicity, and hydrogen-bond donor and acceptor counts summarize interaction potential. For descriptor $k$, standardization is
\begin{equation}
\widetilde d_{ik}=\frac{d_{ik}-\mu_{k,\mathrm{train}}}{\sigma_{k,\mathrm{train}}}.
\end{equation}
The saved training mean and standard deviation are reused unchanged for validation, test and user inputs. Binary fingerprint bits are concatenated with these standardized descriptors in a fixed order. Swapping two descriptor columns would silently change the meaning of the neural input, so feature parity is tested numerically.

The repository also contains variance-filtered and PCA feature artifacts from exploratory feature engineering. The production comparisons described here use the original combined 2,053-dimensional representation; those alternative artifacts must not be assumed to be inputs to the saved models.


\subsection{Multitask Neural Network}
The saved network maps 2,053 inputs through shared layers of widths 512 and 256, followed by a 128-dimensional bottleneck and 12 output logits. Batch normalization and ReLU are used. Dropout is 0.35 in the shared layers and 0.175 at the bottleneck. Sigmoid transforms each logit independently during inference. The implementation uses Adam, weight decay and validation-based early stopping; exact training settings are stored in Member 2's configuration.

\begin{figure}[t]
\centering
\fbox{\parbox{0.90\columnwidth}{\centering
SMILES $\rightarrow$ Morgan bits + scaled descriptors\\[4pt]
$\mathbf{x}\in\mathbb{R}^{2053}$\\[3pt]
Shared network: $512\rightarrow256\rightarrow128$\\[3pt]
$\mathbf{z}\in\mathbb{R}^{128}$\\[5pt]
\textbf{DL:} $\mathbf{z}\rightarrow12$ logits $\rightarrow$ sigmoid\\[3pt]
\textbf{Hybrid:} $[\mathbf{z};\mathbf{d}]\in\mathbb{R}^{133}\rightarrow12$ forests}}
\caption{Shared representation and the two prediction routes. The standalone RF consumes the original 2,053 features directly.}
\label{fig:architecture}
\end{figure}

Let $m_{ij}$ indicate whether molecule $i$ has an observed label for endpoint $j$, and let $w_j$ be the training negative-to-positive ratio. The masked weighted binary loss has the form
\begin{align}
\ell_{ij} &= -w_jy_{ij}\log p_{ij} \\
&\quad -(1-y_{ij})\log(1-p_{ij}),\nonumber\\
\mathcal{L} &= \frac{\sum_{i,j}m_{ij}\ell_{ij}}{\sum_{i,j}m_{ij}}.
\end{align}
Missing labels receive no loss contribution. They must be replaced by finite placeholders before computing elementwise loss and applying the mask. Using a numerically stable loss directly on logits avoids explicit logarithms in code.


\subsection{Optimization and Regularization Details}
The neural configuration uses Adam with initial learning rate $10^{-3}$, weight decay $10^{-4}$ and minibatch size 64. Training is capped at 100 epochs. Early stopping uses validation macro ROC-AUC with patience 15, while the configured learning-rate scheduler monitors validation loss, reduces the rate by a factor of 0.5 after five stagnant epochs, and has a minimum rate of $10^{-5}$. These two controls monitor different quantities: one selects the saved checkpoint and the other adjusts optimization.

ReLU introduces nonlinear transformations of the input features. Batch normalization uses batch statistics during training and stored running statistics during inference. Dropout randomly suppresses activations during training, encouraging reduced dependence on particular hidden features; it is disabled at inference. Weight decay discourages excessively large parameters. These mechanisms address different aspects of optimization and overfitting, and their effectiveness must be judged empirically.

A softmax output would force the 12 assay scores to sum to one, incorrectly treating them as mutually exclusive classes. Independent sigmoid outputs are appropriate because multiple assays can be active simultaneously. For output logit $a_{ij}$,
\begin{equation}
p_{ij}=\frac{1}{1+\exp(-a_{ij})}.
\end{equation}
The masked loss is normalized by the number of observed labels in a minibatch, rather than by all possible entries. This prevents missing entries from diluting the loss. However, it does not enforce equal contributions from endpoints with different observation counts. Positive weighting also changes the effective training objective and does not, by itself, calibrate output probabilities.


\subsection{Phase 1: Classical Machine Learning}
The classical study compares Logistic Regression, Decision Tree and Random Forest separately for each assay. Logistic Regression uses balanced class weights and a maximum of 1,000 optimization iterations. The Decision Tree uses balanced weights and maximum depth 10. These provide a linear reference and a single-tree nonlinear reference before the ensemble comparison. Missing labels are excluded independently for each assay, so the number of observed training molecules varies by endpoint. The saved historical macro ROC-AUC values are 0.7152 for Logistic Regression, 0.6428 for Decision Tree and 0.7541 for Random Forest. Their corresponding average precision values are 0.2949, 0.2091 and 0.3962. These historical metrics are not re-labelled as results from the corrected inference bundle.

\subsection{Standalone and Hybrid Random Forests}
The standalone reference uses 100 trees, maximum depth 12, balanced class weighting and random seed 42 on the original 2,053 features. The hybrid freezes the neural encoder and concatenates its 128-dimensional representation with five standardized descriptors. One Random Forest is trained for each endpoint on observed training labels. The resulting 133-dimensional input tests whether a shared learned representation benefits endpoint-specific trees. Validation and test embeddings are extracted without updating the encoder.


\subsection{Theoretical Comparison of the Model Families}
Logistic Regression learns a linear boundary in the fixed feature space. It provides a useful check on whether predictive structure is captured by approximately additive feature effects. A Decision Tree partitions feature space through sequential tests; it can express interactions but may be sensitive to the particular training sample. Random Forest averages predictions from multiple randomized trees, reducing dependence on one partitioning structure.

The neural model learns intermediate features rather than using only fixed input coordinates. Its shared layers allow every observed endpoint to update a common encoder. The hybrid then changes the classifier applied to that encoder output. Thus, the standalone RF versus hybrid comparison involves both a representation change and a downstream input-size change. It cannot isolate compression alone as a cause of performance differences.

\subsection{Hybrid Training Procedure}
First, a neural checkpoint is selected using validation performance. Second, the encoder is placed in evaluation mode and frozen. Third, embeddings are generated for all partitions without gradient updates and concatenated with the corresponding five scaled descriptors. Finally, each forest is fitted only on training molecules with an observed label for its endpoint. Validation and test labels are excluded from this fitting step. The separate corrected demo bundle follows the same procedure with embeddings regenerated from the actual saved checkpoint.


\section{Evaluation and Results}
ROC-AUC and average precision (called PR-AUC in the repository CSVs) are reported per endpoint and macro-averaged. These answer different questions under imbalance; thresholded accuracy alone is inadequate. Table~\ref{tab:main} uses the historical benchmark in Member 4's \texttt{final\_model\_comparison.csv}. F1 values refer to the saved evaluation protocol, not newly optimized inference thresholds.

\begin{table}[ht]
\caption{Historical team benchmark, macro-average over 12 endpoints}
\label{tab:main}
\centering
\begin{tabular}{lrrr}\toprule
Model & ROC-AUC & Avg. precision & F1\\\midrule
Standalone RF & 0.7541 & 0.3962 & 0.3969\\
Multitask DL & 0.7179 & 0.3376 & 0.3590\\
Hybrid RF & 0.6753 & 0.3009 & 0.2933\\\bottomrule
\end{tabular}
\end{table}

The standalone RF leads this benchmark. The hybrid is successfully implemented but does not improve ranking performance. Compression may remove useful endpoint-specific information; these results alone do not prove that mechanism. The extra descriptors change hybrid ROC-AUC only slightly: the embedding-only ablation reports 0.6738 versus 0.6753 with descriptors.


\subsection{Metric Definitions and Reporting Protocol}
ROC-AUC measures ranking across classification thresholds and can be interpreted as the probability that a randomly selected active example receives a higher score than a randomly selected inactive example, with the usual handling of ties. Average precision summarizes precision at successive increases in recall. The implementation uses average precision, which is not identical to trapezoidal integration of a precision--recall curve.

For true positives $TP$, false positives $FP$, false negatives $FN$ and true negatives $TN$,
\begin{align}
\mathrm{Precision}&=\frac{TP}{TP+FP}, &
\mathrm{Recall}&=\frac{TP}{TP+FN},\\
F_1&=\frac{2TP}{2TP+FP+FN}, &
\mathrm{Balanced\ Accuracy}&=\frac{1}{2}\left(\frac{TP}{TP+FN}+\frac{TN}{TN+FP}\right).
\end{align}
Undefined divisions require an explicit implementation convention. Accuracy alone can be misleading when inactive compounds dominate. Average precision complements ROC-AUC by emphasizing the quality of positive retrieval, and must be interpreted relative to active prevalence.

A macro score is the arithmetic mean of the 12 endpoint scores, giving each assay equal weight. It differs from pooling all molecule--assay predictions into a single micro score. Ranking metrics are threshold-free; F1, recall and confusion matrices depend on the selected threshold. The live demonstration uses a stated default of 0.5. Any optimized threshold must be chosen on validation labels and frozen before test evaluation.

\subsection{Uncertainty and Fair Comparison}
The original classical evaluation includes bootstrap estimates and paired comparisons. Such comparisons require predictions for exactly the same observed test molecules for each endpoint. Their conclusions do not automatically transfer to a separately refitted model bundle. Likewise, the standard deviation across three random seeds measures training variation under a fixed split, not uncertainty over all possible chemical populations. The corrected demo table reports descriptive point estimates; no new claim of statistical significance is made.


\subsection{Separate Three-Seed Experiments}
Member 4 evaluates multitask and single-task learning across seeds 42, 21 and 7. Its summary reports multitask ROC-AUC $0.7472\pm0.0065$ and average precision $0.3286\pm0.0145$, versus single-task ROC-AUC $0.7020\pm0.0076$. The mean ROC-AUC difference is 0.0452. These are experiment means with standard deviations across seeds, not confidence intervals or verified saved-checkpoint metrics. They must not replace the historical row in Table~\ref{tab:main}.

Dropout removal, loss-weight removal and single-task training form the principal neural ablations. The committed summary reports mean active recall 0.6416 for weighted multitask learning and 0.1795 for unweighted learning. This is threshold-dependent and should not be interpreted as universal superiority of weighting at all operating points. Member 4 provides endpoint-level confusion matrices, descriptor error profiles and validation threshold analysis. Error patterns suggest areas for investigation rather than proving biochemical causes.


\subsection{Reading the Ablation Results}
The original-dropout three-seed configuration reports mean ROC-AUC 0.7472 and average precision 0.3286. Removing dropout yields 0.7476 and 0.3207 respectively. The mixed direction of these changes does not support a blanket claim that dropout improves every metric. The weighted-loss configuration has mean recall 0.6416, compared with 0.1795 without weighting at the evaluated operating point. Yet the unweighted model's mean average precision is 0.3489, compared with 0.3286 for the weighted model. This illustrates the distinction between ranking performance and threshold-specific sensitivity.

The single-task comparison asks whether separate networks lose information-sharing benefits. Its lower mean ROC-AUC is consistent with a benefit from shared learning under the tested settings, but the small number of seeds and fixed partition limit generalization. The descriptor ablation shows a small historical hybrid difference, 0.0015 ROC-AUC; this should not be overstated without an uncertainty analysis.

\subsection{Error Analysis Strategy}
False positives and false negatives should be examined per assay rather than collapsed into a single toxicity label. The project provides endpoint metrics, confusion matrices and molecular error inspection. A useful investigation begins with an assay's positive prevalence, then considers score distributions, structural families and descriptor differences between correct and incorrect predictions. These observations can motivate hypotheses about distribution shift or representation limitations.

A low average precision on a rare endpoint can coexist with a moderate ROC-AUC. Conversely, a high recall achieved by flagging many compounds can impose a large false-positive burden. The desired operating point depends on downstream screening priorities and is not established by this project. Structural explanations inferred from error cases remain hypotheses unless supported by additional chemical or experimental evidence.


\section{Inference and Demonstration}
Member 5 validates SMILES with RDKit, reproduces training feature order and uses the committed scaler. The checkpoint is loaded on CPU, switched to evaluation mode and executed with gradients disabled. The same bottleneck feeds the hybrid forests. Standalone RF models are exported from observed training labels using the baseline settings; no fitting occurs in Streamlit.

The interface shows all 12 assay scores for the available models, a comparison chart and downloadable CSV. Invalid molecules and missing artifacts produce explicit messages. The default class threshold is 0.5, which is identified as a default rather than a validation-optimal threshold. Scores are not calibrated probabilities of human toxicity. Integration tests compare new-input features with stored artifacts, compare embeddings against the regenerated checkpoint export and check deterministic prediction output.

\subsection{Member 5 Integration Correction}
Feature parity passed, but the original exported embeddings did not match the saved checkpoint. The trainer retained a mutable state dictionary, so the apparent best weights could change as training continued. It now reloads the on-disk checkpoint before final extraction. Original experiments and artifacts are preserved. A separate demo hybrid is fitted on embeddings regenerated from the saved checkpoint, using only observed training labels. Its test ROC-AUC is 0.7182 with average precision 0.3335. The same bundle reports 0.7318/0.3419 for DL and 0.7498/0.3959 for RF (ROC-AUC/average precision). These integration measurements are not substitutes for the original experiment table. Library versions and checkpoint hash are recorded alongside the demo bundle. No hyperparameters are selected from test scores.

\begin{table}[t]
\caption{Verified inference bundle on the frozen test partition}
\label{tab:verified}\centering
\begin{tabular}{lrr}\toprule
Model & ROC-AUC & Average precision\\\midrule
Standalone RF & \textbf{0.7498} & \textbf{0.3959}\\
Saved multitask checkpoint & 0.7318 & 0.3419\\
Rebuilt hybrid RF & 0.7182 & 0.3335\\\bottomrule
\end{tabular}
\end{table}

\subsection{Endpoint-Level Results for the Verified Bundle}
Table~\ref{tab:endpoints} exposes assay variation that is hidden by a macro-average. All entries come from the corrected bundle's committed test metrics. Endpoint winners can differ, and a stronger aggregate score does not imply dominance on every assay.
\begin{table}[htbp]\centering\small
\caption{Verified test ROC-AUC and average precision (AP) by endpoint}
\label{tab:endpoints}
\begin{tabular}{lrrrrrr}\toprule
& \multicolumn{2}{c}{Random Forest} & \multicolumn{2}{c}{Multitask DL} & \multicolumn{2}{c}{Hybrid}\\
Endpoint & ROC & AP & ROC & AP & ROC & AP\\\midrule
NR-AR & 0.6922 & 0.4848 & 0.7505 & 0.4595 & 0.7449 & 0.4383\\
NR-AR-LBD & 0.8177 & 0.4535 & 0.7731 & 0.4151 & 0.8130 & 0.4502\\
NR-AhR & 0.8207 & 0.5245 & 0.7803 & 0.4249 & 0.7621 & 0.3982\\
NR-Aromatase & 0.7811 & 0.4623 & 0.7454 & 0.3099 & 0.7398 & 0.2929\\
NR-ER & 0.6342 & 0.3529 & 0.6337 & 0.3347 & 0.6270 & 0.3196\\
NR-ER-LBD & 0.7221 & 0.3825 & 0.7273 & 0.3290 & 0.7263 & 0.3406\\
NR-PPAR-gamma & 0.7752 & 0.2371 & 0.7750 & 0.2269 & 0.6984 & 0.2258\\
SR-ARE & 0.7267 & 0.5203 & 0.6968 & 0.4050 & 0.7044 & 0.4409\\
SR-ATAD5 & 0.7513 & 0.1607 & 0.6412 & 0.0954 & 0.6111 & 0.0997\\
SR-HSE & 0.7296 & 0.2681 & 0.7369 & 0.2319 & 0.7003 & 0.1994\\
SR-MMP & 0.8239 & 0.6102 & 0.8122 & 0.5974 & 0.8050 & 0.5314\\
SR-p53 & 0.7229 & 0.2944 & 0.7095 & 0.2732 & 0.6855 & 0.2649\\
\bottomrule\end{tabular}\end{table}

\subsection{Interpretation of the Verified Results}
In Table~\ref{tab:verified}, RF exceeds DL by 0.0180 ROC-AUC and 0.0540 average precision. The rebuilt hybrid trails DL by 0.0136 ROC-AUC. These descriptive differences do not establish statistical significance. The corrected bundle demonstrates executable consistency while retaining the same scientific conclusion: the tested hybrid does not outperform the strongest baseline. Comparisons across historical and corrected runs also reflect artifact and software differences and should not be interpreted as a controlled ablation.


\subsection{Unified Dashboard and User Workflow}
The application presents the project as one system through six sections: Overview, Data and Molecular Features, Model Architectures, Experiments and Results, Predict a Molecule, and Report and Presentation. The overview summarizes the dataset and verified model comparison. The data section explains label coverage and feature construction. The architecture section exposes classical models, the neural network, its training history and the hybrid route. The experiment section distinguishes historical results, the verified bundle and ablation studies.

For inference, the user enters a SMILES string or chooses a preset. Invalid or empty structures are rejected before feature extraction. A valid input is canonicalized, featurized and evaluated by the saved models. The output contains 36 scores: 12 endpoints for each of three model families. The interface also renders the molecular structure and exports a CSV containing endpoint, model, activity score, predicted class, threshold and canonical SMILES. Results remain associated with the molecule that was actually evaluated.

\subsection{Verification Performed}
The inference tests cover invalid-input rejection, agreement between regenerated features and stored artifacts, agreement between regenerated embeddings and the corrected checkpoint export, and deterministic prediction output. The existing neural-model test suite also passes. Dashboard checks exercise all six navigation sections and verify a complete live prediction with 36 scores. These checks support functional integration; they do not establish calibration, clinical usefulness or robustness to every possible molecule.

The corrected bundle records the neural checkpoint's SHA-256 digest and software versions. Inference checks the recorded digest before loading the model, so a changed checkpoint cannot silently reuse incompatible downstream artifacts. The exported forest files are intended to be loaded only from the trusted project bundle. Rebuilding after a checkpoint change must regenerate both matching embeddings and hybrid forests.


\section{Limitations and Future Work}
The dataset is small and imbalanced, and the single frozen scaffold partition does not capture all prospective distribution shifts. Three training seeds quantify only part of experimental variation. Some original documentation mixes production results and separate ablation summaries; this report keeps their provenance distinct. Hybrid improvement has not been demonstrated under the tested configuration. Neither feature associations nor error profiles establish causal toxicological mechanisms. Future work may assess calibration on validation data, alternative bottleneck sizes and repeated scaffold partitions while maintaining an untouched evaluation protocol.

\section{Team Contributions and Reproducibility}
Member 1 owns the split contract, masked data loader and literature table. Member 2 owns the multitask network, training and embeddings. Member 3 owns the hybrid forests. Member 4 owns evaluation and ablations. Member 5 owns inference integration, the demonstration, submission assembly and operational documentation. Individual attribution and commit evidence are maintained in the accompanying contribution record.

Install root requirements and launch \texttt{streamlit run app.py}. If needed, rebuild the inference bundle using the \texttt{prepare\_models} module documented in the Member 5 guide. Run the Member 5 test suite with unittest discovery in its tests directory. Exact commands and demonstration steps are in the adjacent setup and reproduction guide. Reported model metrics come from committed team experiment artifacts; Member 5 does not claim to have rerun all training experiments.


\subsection{Reproduction Sequence and Artifact Ownership}
A complete reproduction begins with the frozen data artifacts and their validation summary. Classical models consume the combined feature matrix. Neural training consumes the same matrix with label masks and training-derived positive weights. The saved checkpoint must then be restored before any embedding export. Hybrid fitting uses those embeddings and the matching descriptor rows. Evaluation is performed with unchanged split membership, followed by integration parity checks and dashboard prediction tests.

The project stores model configuration, training history, checkpoint files, feature/scaler artifacts, endpoint predictions and aggregate metrics separately. This separation makes it possible to inspect the origin of a result and rebuild only the stage that changed. Model files alone are insufficient for reproduction: feature order, scaler parameters, endpoint order and package versions are also part of the prediction contract.

For submission, the report, slides and dashboard should cite the same experiment family for each number. The contribution record should be completed with actual names, roll numbers and commit evidence. The repository provides implementation and result evidence for both phases, while individual understanding must be demonstrated through explanation of the methods, experiments and limitations.


\section{Conclusion}
The project establishes a reproducible comparison of fingerprint-based machine learning, multitask deep learning and hybrid representation learning. Random Forest remains the strongest measured method, while the separate seed study supports investigating shared learning across assays. The inference audit shows why saved weights, preprocessing and downstream features must be verified together. The resulting demonstration serves all 12 endpoints from a documented model bundle, with transparent reporting of both negative results and unresolved limitations.

\begin{thebibliography}{4}
\bibitem{ramsundar} B. Ramsundar et al., ``Massively Multitask Networks for Drug Discovery,'' arXiv:1502.02072, 2015.
\bibitem{moleculenet} Z. Wu et al., ``MoleculeNet: a benchmark for molecular machine learning,'' Chemical Science, vol. 9, pp. 513--530, 2018.
\bibitem{morgan} D. Rogers and M. Hahn, ``Extended-Connectivity Fingerprints,'' Journal of Chemical Information and Modeling, vol. 50, pp. 742--754, 2010.
\bibitem{breiman} L. Breiman, ``Random Forests,'' Machine Learning, vol. 45, pp. 5--32, 2001.
\end{thebibliography}
\end{document}
